\documentclass[10pt,a4paper]{article}
\usepackage[utf8]{inputenc}
\usepackage[T1]{fontenc}
\usepackage[spanish,es-noshorthands,es-tabla]{babel}
\usepackage{lmodern}
\usepackage{microtype}
\usepackage[margin=1.15cm,top=1.1cm,bottom=1.1cm]{geometry}
\usepackage{multicol}
\usepackage{booktabs}
\usepackage{tabularx}
\usepackage{xcolor}
\usepackage{tikz}
\usetikzlibrary{arrows.meta,positioning,babel}
\usepackage{hyperref}

% Paleta de colores sobria y académica
\definecolor{azul}{HTML}{14334B}
\definecolor{celeste}{HTML}{EAF3F8}
\definecolor{acento}{HTML}{006699}
\definecolor{gris}{HTML}{4A5560}

\hypersetup{
  colorlinks=true,
  urlcolor=acento,
  linkcolor=acento
}

\pagestyle{empty}
\setlength{\parindent}{0pt}
\setlength{\parskip}{3pt}
\setlength{\columnsep}{0.6cm}

\newcommand{\seccion}[1]{\par\vspace{4pt}{\color{azul}\bfseries\normalsize #1}\par\vspace{1pt}}

\begin{document}

% Encabezado institucional compacto y elegante
{\color{azul}\bfseries\large Solución RAG para el asistente de normativa de pregrado (Ingeniería UdeC)}\par
\vspace{1pt}
{\color{acento}\bfseries\footnotesize Entregable 2: recuperación densa y prompt estructurado sobre Qwen3-4B}\par
\vspace{2pt}
{\scriptsize \textbf{Generative Artificial Intelligence (580694)} \textbar{} Segundo semestre 2026 \textbar{} \textbf{Equipo:} Álvaro Contreras y Pablo Cortés}\par
\vspace{1pt}
{\scriptsize \textbf{Repositorio:} \href{https://github.com/Pacortes2021/GENERATIVE-ARTIFICIAL-INTELLIGENCE-DELIVERABLE-1}{github.com/Pacortes2021/GENERATIVE-ARTIFICIAL-INTELLIGENCE-DELIVERABLE-1} \hfill \textbf{Video (3 min):} \href{https://drive.google.com/file/d/1eBhjtr9q2l0yT9FnfQ_rWjtIsbdhvvY0/view?usp=drive_link}{[Enlace de video]}}\par
\vspace{2pt}\hrule height 0.8pt\vspace{3pt}

\begin{multicols}{2}

\seccion{1. Continuidad con E1 y selección de modelo}
Este entregable continúa la misma tarea del Entregable~1: responder preguntas en español sobre la normativa de pregrado de la FI-UdeC a partir de un corpus cerrado de tres documentos (Reglamento General, Reglamento Interno FI y Calendario 2026). La corrección es estricta (dato exacto más cita oficial del artículo; abstención si no existe o la premisa es falsa), evaluada manualmente sobre las mismas 50 preguntas.

De los tres candidatos de E1, ratificamos \textbf{Qwen/Qwen3-4B} (revisión \texttt{1cfa9a720891}, Apache~2.0) por ser el menor en parámetros. La falla diagnosticada fue \emph{ausencia paramétrica} (el modelo inventa al no tener el corpus en sus pesos). La solución inyecta la evidencia mediante RAG. La inferencia corre en Google Colab (GPU NVIDIA Tesla T4, 15\,GB VRAM) en cuantización NF4 de 4 bits, cómputo en \texttt{bfloat16} y decodificación greedy (\texttt{temperature=0}, sin modo thinking).

\seccion{2. Arquitectura del sistema y flujo comparativo}
\begin{center}
\begin{tikzpicture}[
  >=Latex, node distance=2.5mm,
  box/.style={draw=azul,rounded corners=2pt,fill=celeste,align=center,font=\tiny,minimum height=5.2mm,inner sep=2.5pt},
  arr/.style={->,draw=azul,line width=.5pt},
  darr/.style={->,draw=acento,dashed,line width=.5pt}
]
% Rama Baseline
\node[box] (q) {Pregunta};
\node[box,right=7mm of q] (dir) {Prompt directo\\(sin documentos)};
\node[box,right=5mm of dir] (m1) {Qwen3-4B\\(NF4)};
\node[box,right=5mm of m1] (out1) {Respuesta\\Baseline};

% Rama RAG
\node[box,below=5.5mm of q] (idx) {Corpus: 3 PDFs\\(197 unid. / 201 frag.)};
\node[box,right=4mm of idx] (e5) {Índice E5\\multilingual};
\node[box,below=5.5mm of m1] (m2) {Qwen3-4B\\prompt estruct.};
\node[box,right=5mm of m2] (out2) {Respuesta RAG\\(dato + cita)};

\draw[arr] (q) -- (dir);
\draw[arr] (dir) -- (m1);
\draw[arr] (m1) -- (out1);

\draw[arr] (idx) -- (e5);
\draw[darr] (q.south) |- (e5.west);
\draw[arr] (e5) -- (m2);
\draw[arr] (m2) -- (out2);
\end{tikzpicture}
\end{center}
\vspace{-4pt}
Los PDFs se segmentaron en \textbf{197 unidades normativas} (201 fragmentos de búsqueda). Con \texttt{multilingual-e5-small} se recuperan los fragmentos afines, expandiéndolos a la unidad íntegra y siguiendo remisiones de un salto para evitar artículos truncados. El prompt estructurado exige dato suficiente, condiciones y cita. El baseline recibe la pregunta sin documentos con la instrucción directa de E1.

\seccion{3. Estrategias alternativas evaluadas}
Sobre las 50 preguntas de desarrollo se contrastaron tres formulaciones con la misma recuperación E5:
\begin{center}
\scriptsize
\begin{tabularx}{\linewidth}{@{}Xrcc@{}}
\toprule
\textbf{Variante} & \textbf{Aciertos} & \textbf{Exactitud} & \textbf{Decisión} \\
\midrule
RAG simple (instrucción breve) & 39 / 50 & 78\% & Alternativa \\
\textbf{RAG estructurado} & \textbf{40 / 50} & \textbf{80\%} & \textbf{Elegida} \\
RAG con ejemplos (few-shot) & 36 / 50 & 72\% & Descartada \\
\bottomrule
\end{tabularx}
\end{center}
\vspace{-3pt}
Se seleccionó RAG estructurado por su formato explícito y control de alucinaciones. Few-shot se descartó porque los ejemplos ficticios del prompt contaminaron citas normativas.

\seccion{4. Resultados frente a la línea base}
En evaluación emparejada (misma GPU T4, pesos, NF4, parámetros y semilla por pregunta):
\begin{center}
\scriptsize
\begin{tabularx}{\linewidth}{@{}Xcc@{}}
\toprule
\textbf{Categoría (10 c/u)} & \textbf{Baseline directo} & \textbf{RAG estructurado} \\
\midrule
Factual & 0 / 10 & 9 / 10 \\
Numérica & 0 / 10 & 7 / 10 \\
Condicional & 0 / 10 & 6 / 10 \\
Cruce de fuentes & 0 / 10 & 8 / 10 \\
Premisa falsa / abstención & 1 / 10 & 10 / 10 \\
\midrule
\textbf{Global ($n=50$)} & \textbf{1 / 50 (2\%)} & \textbf{40 / 50 (80\%)} \\
\bottomrule
\end{tabularx}
\end{center}
\vspace{-3pt}
RAG mejoró 78 puntos porcentuales (+39 preguntas netas). Acertó en 10/10 casos de abstención/premisa falsa, eliminando la invención de reglamentos inexistentes. El 4\% histórico de E1 se mantiene separado por provenir de otra configuración de inferencia. La latencia promedio fue 16,1\,s (RAG) vs. 7,2\,s (baseline), con recuperación en 0,11\,s.

\seccion{5. Fallos reales y límites conocidos}
Los 10 fallos de RAG estructurado son de síntesis, omisión de condiciones o recuperación, no de invención:
\begin{itemize}
  \item \textbf{P25 (fallo de síntesis):} Ante la duda de qué asignaturas no se pueden eliminar al modificar inscripción, respondió solo \guillemotleft{}prioridades a) y b)\guillemotright{}. Los Arts. 7 y 9 de RI-FI estaban íntegros en el contexto, pero el generador no desarrolló que correspondían a asignaturas atrasadas y reprobadas.
  \item \textbf{P38 (fallo de recuperación):} E5 no recuperó el Art. 2 relevante y trajo el Art. 6 (cupos), provocando una cita errónea.
\end{itemize}
Las 50 preguntas guiaron el ajuste del recuperador: este 80\% no mide generalización a consultas no vistas.

\seccion{6. Reproducibilidad y video de ejecución}
El cuaderno \texttt{Deliverable2\_Sistema\_RAG\_Colab.ipynb} es autocontenido en Colab: carga E5, el corpus procesado y Qwen3-4B, ejecutando ambas condiciones en vivo sobre la misma entrada. El video adjunto (máx. 3\,min) muestra el sistema corriendo de extremo a extremo con P25 y la auditoría de veredictos. El repositorio preserva hashes SHA-256 de prompts, corpus, los 100 JSON de respuesta y los veredictos auditados caso por caso.

\end{multicols}
\end{document}
